\section{Uniform analytic lifts and genuine physical continuation tubes}
\label{sec:pt-tubes}
All radii, disks, collision labels, physical clocks and the original
section are unchanged.  The constants in this section are deliberately
conservative; their purpose is a positive radius independent of the
number of actual collisions.  Unlike a bound for each fixed derivative,
the first theorem provides one common domain of analytic continuation.
A separate geometric argument determines when this continuation is a
true billiard trajectory.  The latter step is not inferred from analyticity.

Let $a=47/6$, $\lambda=50/47$, $v_i=\lambda^{\min(i,m-i)}$, and set
\begin{align}
 K&=2(1+\lambda)(4a+16a^2),
 &r_*&=\min\{[8(1+a)]^{-1},(8K)^{-1}\},\notag\\
 \eta_*&=\frac{r_*}{8(1+\lambda a)},
 &P_*&=2(1+\lambda a),\label{eq:pt-analytic-constants}\\
 C_{\rm inc}&=P_*[1+2a(1+\lambda)],
 &C_{\rm clr}&=R_+P_*,\notag\\
 \delta_*&=\min\{\eta_*,(2C_{\rm inc})^{-1},
                              (2C_{\rm clr})^{-1}\}.\label{eq:pt-tube-constant}
\end{align}
The symbols $r_*,\eta_*,\delta_*$ in this section are auxiliary
endpoint-domain constants, not the original microscopic width $r_B$.
For complex vectors put $\|z\|_v=\sum_{i=0}^m v_i|z_i|$.

\begin{lemma}[Uniform complex neighborhood of the reflection equations]
\label{lem:pt-complex}
At every actual regular reference trajectory $\alpha^0$, the reflection
vector $E$ of \eqref{eq:gr-reflection} has a holomorphic continuation to
$\|\alpha-\alpha^0\|_v<r_*$.  On that ball,
\[
 \|D^2E(\alpha)[z,w]\|_v\le K\|z\|_v\|w\|_v.
\]
The output norm uses only interior indices.  Both complexified flight
coordinates stay nonzero, uniformly in collision count and incidence.
\end{lemma}
\begin{proof}
Write the two complexifications of a flight as
\[
 X_i^+=(O_{i+1}-O_i)_x+i(O_{i+1}-O_i)_y
             +R(e^{i\alpha_{i+1}}-e^{i\alpha_i}),
\]
and let $X_i^-$ be the same expression with $i$ replaced by $-i$.
For complex angles $X_i^-$ is not defined by complex conjugation.
At the real reference point $|X_i^+|=|X_i^-|=\tau_i\ge\tau_-$.
For $\|\alpha-\alpha^0\|_v\le r_*$,
\[
 |X_i^\pm(\alpha)-X_i^\pm(\alpha^0)|
 \le R_+e^{r_*}(|\alpha_i-\alpha_i^0|
                      +|\alpha_{i+1}-\alpha_{i+1}^0|)
 \le2R_+r_*<\tau_-/4.
\]
Here $r_*<1/8$, $e^{r_*}<2$, and $2ar_*<1/4$.
Choose the logarithms on these nonvanishing disks and put
$\vartheta_i=(\log X_i^+-\log X_i^-)/(2i)$, with the additive
real lift fixed at the reference point.  This agrees with the real
flight angle.  Each first or second derivative of $X_i^\pm$ has
bound $2R_+\prod_h(|z_i^h|+|z_{i+1}^h|)$, and
$|X_i^\pm|\ge\tau_-/2$.  The two logarithmic derivative terms
are consequently bounded by $4a$ and $16a^2$.
The weighted flight summation in Lemma~\ref{lem:gr-angle-tensors}
then gives $2(1+\lambda)(4a+16a^2)$, without a factor $m$.
\end{proof}

\begin{theorem}[A collision-uniform analytic endpoint domain]
\label{thm:pt-analytic}
Fix any actual regular $m$-flight trajectory, $R\in[R_-,R_+]$,
and its designated disk centers.  For every complex endpoint increment
$h=(h_0,h_m)$ with $s=|h_0|+|h_m|<\eta_*$ there is a unique
solution of the same locally lifted reflection equations in the
interior ball $\|\alpha_I-\alpha_I^0\|_v\le r_*/2$.
It depends holomorphically on $h$, is real for real $h$, and satisfies
\begin{equation}\label{eq:pt-analytic-localization}
 \sum_{i=0}^m v_i|\alpha_i(h)-\alpha_i^0|\le P_*s.
\end{equation}
Uniqueness is asserted in this ball, not among arbitrary distant
solutions of the word equations.  A simultaneous visited-grazing
limit of reference configurations has the same formal statement.
No physical passage through grazing is asserted.
\end{theorem}
\begin{proof}
For $m=1$ the assertion is just the two endpoints.  Otherwise let
$B_0=D_IE(\alpha^0)$ and let $Q_0h$ be its boundary forcing.
Theorem~\ref{thm:gr-resolvent} gives
$\|B_0^{-1}\|_{v\to v}\le1$ also over complex vectors, and
$\|Q_0h\|_v\le\lambda a s$.  Write
\[
 E(\alpha^0+(h,u))-E(\alpha^0)=B_0u+Q_0h+N(h,u).
\]
On the stated ball $s+\|u\|_v<r_*$, the preceding lemma and
integration along the complex line segment give
$\|N(h,u)\|_v\le K(s+\|u\|_v)^2/2$ and
$\|D_uN\|\le K(s+\|u\|_v)$.
The map $u\mapsto-B_0^{-1}(Q_0h+N(h,u))$ sends
$\|u\|_v\le r_*/2$ into itself: its forcing is at most $r_*/8$
and its nonlinear part is at most $r_*/16$.  Its Lipschitz constant
is at most $Kr_*\le1/8$.  The successive iterates from zero thus
converge geometrically, uniformly on every smaller endpoint domain,
to the unique fixed point.  The iterates are holomorphic, so their
locally uniform limit is holomorphic.  Real coefficients and uniqueness
give reality for real $h$.
For the solution, $s+\|u\|_v\le r_*$ yields
\[
 \|u\|_v\le\lambda a s+\tfrac1{16}(s+\|u\|_v),
 \qquad
 s+\|u\|_v\le\tfrac{16}{15}(1+\lambda a)s\le P_*s.
\]
This proves \eqref{eq:pt-analytic-localization}.
At a fixed-word grazing limit the real limiting flight lengths are
still at least $\tau_-$ and the limiting reflection inverse retains
its weighted bound.  The identical contraction proof applies there.
It concerns reflection equations only until physical admissibility
is verified separately below.
\end{proof}

\begin{theorem}[Physical tubes for the original depth-dependent guard]
\label{thm:pt-physical}
Put $r=q_0^{1/4}$, with the unchanged $q_0=47/53$.  Suppose an actual
regular reference orbit satisfies, for $0<\varepsilon\le1/4$,
\begin{equation}\label{eq:pt-guarded-center}
 c_i^0\ge\varepsilon r^{d_i}\quad(0\le i\le m),\qquad
 \Delta_i^0\ge\varepsilon r^{d'_i}\quad(0\le i<m).
\end{equation}
For every real $h$ with $|h_0|+|h_m|<\delta_*\varepsilon$, the
analytic lift is an actual regular trajectory with exactly the same
ordered collisions and displacement.  Every designated incidence and
every least third-disk clearance obeys
\begin{equation}\label{eq:pt-half-margin}
 c_i(h)\ge\tfrac12\varepsilon r^{d_i},\qquad
 \Delta_i(h)\ge\tfrac12\varepsilon r^{d'_i}.
\end{equation}
The physical clock is the sum of the true lengths.  The original
section, visit count and terminal test are evaluated at the resulting
physical orbit; they are not declared constant on the tube.
Thus any Borel exact-record selector can be retained pointwise without
adding a section-distance condition to the guard.
\end{theorem}
\begin{proof}
The norm bound in \eqref{eq:pt-analytic-localization} gives
$|\alpha_i(h)-\alpha_i^0|\le P_*s\lambda^{-d_i}$.
For real angles the corresponding endpoint displacement is at most
$R_+P_*s\lambda^{-d_i}$.  The Hausdorff distance between the two
closed flight segments is at most
$C_{\rm clr}s\lambda^{-d'_i}$.
Distance to every third disk is one-Lipschitz for this distance;
taking the infimum over all centers except exactly the two incident
centers preserves it.  This controls every unvisited obstacle at once,
including switches of the minimizing obstacle or closest segment point.

For a nonzero real flight, comparison of normalized vectors gives
$|X/|X|-X^0/|X^0||\le2|X-X^0|/|X^0|$.
Since adjacent depth weights have ratio at most $\lambda$, each
signed departure normal component, and each negative arrival normal
component, changes by at most
$C_{\rm inc}s\lambda^{-d_i}$.  The same bound holds for the
signed tangential component.  Initially both physical normal
components at an interior reflection equal $c_i^0$.
The exact rational inequality
\[
 \lambda^4q_0=\frac{50^4}{47^3\,53}>1
\]
implies $\lambda^{-d}\le r^d$.  The choice of $\delta_*$ therefore
proves both half-margins, first for signed normal components and
segment clearances of the formal real lift.

Their strict positivity makes that lift physical.  A positive outgoing
normal component keeps a segment outside its initial disk; a negative
incoming normal component keeps it outside its final disk until the
stated endpoint.  All other disks have strictly positive clearance.
Thus each segment is the true next free flight, not a flight with a
hidden earlier collision.  The fixed reflection equations, together
with these normal signs, give specular reflection at every interior
endpoint.  Concatenation proves the actual ordered word and unchanged
displacement, with no deleted collision.  Its unit-speed duration is
its geometric length.  Section membership can jump as its boundary is
crossed, so the original selector is retained rather than differentiated
or frozen.  This makes no assertion about the source outside the tube.
\end{proof}

The guard margins in \eqref{eq:pt-guarded-center} already occur in the
original product \eqref{eq:v76-guard}; they are not a new conditioning
event for the target probability.  The theorem establishes an explicit
domain property of that existing geometric class.  Points failing the
margins, and all one-sided singular traces, remain present in the full
source problem.  No claim that their normalized mass tends to zero is
made by constructing these tubes.
